\section{Limitations}
\label{sec:limitations}

The DRM Transformer is a research prototype. The current baseline is small,
approximately 3.5 million parameters, and was trained on few tokens compared
with modern language models. Therefore, perplexity results should not be
directly compared with production models.

Topological validation depends on evaluation choices: number of Voronoi seeds,
homology subsample size, long-bar ratio, PCA projection, density filtering, and
number of restarts. The strict criterion $f_{T^2}\geq 0.60$ reduces false
positives, but remains an experimental convention. In the current results, the
trained model exhibits a near-stable toroidal signature, but does not reach that
strict threshold.

Toroidal losses are inductive: they help induce $T^2$ in dimension four, but
they may also compete with the language loss. The fact that controls without
toroidal regularization still show some near-$T^2$ signal suggests that the
architecture also contributes to the structure, but separating learned
topology, induced topology, and TDA artifacts requires further ablations.

In addition, the local geometry remains thick in 4D across several runs. The
toroidal signature appears more clearly under PCA3 homology with density
filtering, and not as a fully stable signature in the raw 4D embedding. At
larger model scales, geometric loss weights will likely need to be rescaled.

Finally, the connection with relativistic physics is formal and analogical. The
Lorentz factor is used as mathematical inspiration for scale control; the model
does not claim to implement a new physical theory.

Similarly, the discussion of open geometry, negotiation, and AI safety is a
conceptual motivation, not an operational guarantee. The DRM Transformer does
not prove that an AI will be safe because it has closed topology, nor does it
replace behavioral evaluation, governance, interpretability, or external
control. The contribution here is a measurable hypothesis: if some epistemic and
semantic conflicts appear as geometric structures, then they can be compared
through ablations, persistent homology, and initialization controls.
